\section{An ordinary adjoint on the projective base}
\label{sec:adjoint}

The induction will produce a morphism to a projective variety for which an
ordinary adjoint is pulled back from the base.  To apply the projective
numerical-semiampleness theorem, we must realize the line downstairs as an
ordinary klt adjoint.  The following proposition provides this realization.
The total space need not be projective.

\begin{proposition}[An adjoint on the base]
\label{adj:base}
Let $(Y,D)$ be an effective klt pair with rational boundary on a normal
compact K\"ahler space.  Let $f:Y\to Z$ be a surjective morphism with connected
fibers to a normal projective variety, and suppose that
$\dim Y>\dim Z$.  If $H$ is a rational Cartier line bundle on $Z$ and
\begin{equation}
\label{adj:initial}
                         K_Y+D\qsim f^*H,
\end{equation}
then there is an effective rational divisor $D_Z$ such that $(Z,D_Z)$ is klt
and
\[
                         H\qsim K_Z+D_Z.
\]
\end{proposition}

Ambro proves the corresponding statement for projective total spaces
\cite[Theorem~0.2]{Ambro}.  We follow the same separation into a discriminant
and a Hodge-theoretic moduli line.  The two points requiring attention here
are the real, possibly irrational polarization of a K\"ahler family and the
construction of its cyclic cover without a meromorphic frame of $K_Y$.
The period results recalled next address the first point; we give the cover
and the required infinitesimal calculation explicitly.

\subsection{The period input and the discriminant}

For a pure variation of Hodge structures on a smooth open set, the
\emph{line period map} of a rank-one highest Hodge piece is the map to the
projective space of the flat vector space that records this line in a local
flat trivialization.  Its generic differential rank is independent of the
trivialization.  The full period map records the entire Hodge filtration.
For quasi-unipotent boundary monodromy, the \emph{parabolic extension} of the
line is the rational line bundle obtained by making the monodromy unipotent
on local finite covers, extending the highest Hodge piece, and descending
with its rational boundary weights.

We use the following precise consequence of the period results in
\cite[Lemmas~4.1 and~4.3 and Proposition~4.5]{LI}.

\begin{lemma}[Period quotient input]
\label{adj:period-input}
Let $S$ be a smooth projective variety and let $S^\circ\subset S$ have simple
normal crossing complement.  Let $\mathbb V$ be a real-polarizable pure
variation on $S^\circ$ with an integral lattice and quasi-unipotent local
monodromy.  Suppose that a complex direct summand of $\mathbb V_{\C}$, as a
variation of Hodge structures, has a rank-one highest Hodge piece, with parabolic extension $L$.  Then $L$ is nef,
and its numerical dimension is the generic rank of its line period map.

After modification $\tau:S'\to S$, there are a smooth projective variety
$Q$, a surjective morphism $p:S'\to Q$ with connected fibers, and a nef
rational line bundle $P$ on $Q$ such that
\begin{equation}
\label{adj:period-descent}
                         \tau^*L\qsim p^*P.
\end{equation}
Moreover, $\dim Q$ is at most the generic rank of the full period map of
$\mathbb V$.  If $Q$ is a point, $\tau^*L$ is rationally trivial.
\end{lemma}

Here the identity is in $\Pic(S')\otimes\Q$, including the boundary.
For clarity, the dimension assertion follows because the descended,
generically immersive period map on $Q$ belongs to an adjoint variation
obtained by tensor constructions from $\mathbb V$.
The rational adjoint reduction in \cite[Lemma~4.3]{LI} retains entire
rational simple factors, so its monodromy is discrete even when the
original real polarization is irrational.  It also kills only finite
scalar monodromy after taking a power.  Thus
\eqref{adj:period-descent} does not discard an arbitrary flat twist.
We shall need exactly this actual line identity.

Choose a projective resolution $b:S\to Z$, and normalize the main
component of $Y\times_Z S$.  Choose its log resolution functorially
with respect to local isomorphisms preserving the marked boundary;
over the open where $b$ is an isomorphism this is a functorial log
resolution of $(Y,D)$.  Denote the resulting smooth compact K\"ahler
space by $X$, with maps $\pi:X\to Y$ and $g:X\to S$.
Such a resolution is available in the complex analytic category;
see \cite[Theorem~1.1 and the preceding analytic remark]{BierstoneMilman2008}.
The functorial choice will matter when we lift local flows.
Additional resolutions used to compute thresholds and fiber integrals
are auxiliary and do not replace the family whose cohomology we use.
Define the crepant rational divisor $D_X$ by
\[
              K_X+D_X=\pi^*(K_Y+D).
\]
Its exceptional coefficients may be negative; all its coefficients are
strictly less than one.  For a prime divisor $Q_0$ on $S$, let
\[
 t_{Q_0}=\sup\{c\in\R:(X,D_X+c g^*Q_0)
             \text{ is log canonical over the general point of }Q_0\}.
\]
The definition is unchanged on a higher crepant resolution.  It gives a
positive rational number, computed by a log resolution as
\begin{equation}
\label{adj:threshold}
             t_{Q_0}=\min_i\frac{1-u_i}{a_i},
\end{equation}
where $g^*Q_0=\sum_i a_i E_i$ over its general point and $u_i$ is the
coefficient of $E_i$ in the crepant boundary there.  Define
\begin{equation}
\label{adj:discriminant}
 \Delta_S=\sum_{Q_0}(1-t_{Q_0})Q_0,
 \qquad L_S=b^*H-K_S-\Delta_S.
\end{equation}
Only finitely many terms of $\Delta_S$ are nonzero: off a suitable proper
analytic subset the resolved pair is relatively simple normal crossing,
the map is smooth, and the threshold of a base prime is one.  We may choose
$S$ with a simple normal crossing divisor containing this exceptional set
and the support of $\Delta_S$.

The coefficients of $\Delta_S$ are strictly less than one.  To prove
Proposition~\ref{adj:base}, we will show that $L_S$ has an effective rational
representative whose addition to $\Delta_S$ is sub-klt.  The relevant
positivity comes from the following root construction.

\subsection{The root cover and its highest Hodge line}

\begin{lemma}[The root eigenline]
\label{adj:root}
In the setting of Proposition~\ref{adj:base}, there is a dense Zariski open
$S^\circ\subset S$ and a smooth proper family $h:V^\circ\to S^\circ$ obtained
from a cyclic cover and resolution, such that $R^l h_*\R$, where
$l=\dim Y-\dim Z$, is real-polarizable and has an integral lattice.
One character summand of its complexification has a rank-one highest
Hodge piece.  On $S^\circ$, this line agrees, as a rational line bundle,
with $b^*H-K_S$.
\end{lemma}

\begin{proof}
Choose an integer $m>0$ clearing all divisors and rational line identities.
A meromorphic section of the line $m b^*H$ exists because $S$ is projective.
Its pullback, under the identity
\[
                m(K_X+D_X)\simeq g^*(m b^*H),
\]
and division by the canonical meromorphic section of $mD_X$ give a
meromorphic section $\sigma$ of $K_X^{\otimes m}$.
On the open where $\sigma$ has neither zeros nor poles, take its $m$th
roots in the fibers of $K_X$.  In local canonical frames this is the
equation $z^m=\sigma$; on overlaps the roots transform by the transition
functions of $K_X$.  The local covers therefore glue without a
meromorphic trivialization of $K_X$, and without choosing a root of $H$.

The Grauert--Remmert extension theorem for finite analytic covers
\cite[Exposé~XII, Proposition~5.3 and Theorem~5.4]{SGA1} extends it
normally across the divisor of $\sigma$, giving a finite cyclic cover
of $X$, possibly disconnected.  It is compact K\"ahler.
Take a resolution functorial for the cyclic action and for local
isomorphisms of the pair, as in \cite[Section~1.1 and Lemma~1.1]{Ambro}.
Thus both resolutions are functorial.  A local flow preserving the
original pair lifts through them once its lift to the root cover has
been chosen; we will construct that lift explicitly below.
The resulting compact manifold $V$ is K\"ahler.  After deleting from $S$ the zeros and
poles of the chosen base section and the degeneration locus, the map
$h:V^\circ\to S^\circ$ is smooth and proper.  We take $S^\circ$
inside the isomorphism locus of $b$.  Shrinking it further, the original
fibers of $Y\to Z$ are normal effective klt pairs and the chosen
resolutions restrict to resolutions of those pairs.

Average a global K\"ahler class on $V$ under the cyclic group.  Its
restriction is a flat real section of $R^2h_*\R$ and polarizes the
primitive pieces.  Their Lefschetz decomposition supplies a flat real
polarization of the full degree-$l$ cohomology.  We retain the full
cohomology local system, with lattice $R^lh_*\Z$ modulo torsion; an
irrational primitive summand is not required to carry its own lattice.
The finite group action preserves this variation and its polarization.

Locally over the base, the tautological root, divided by a local base
volume form, gives a relative meromorphic top form $\tau$ on the cover.
It belongs to a fixed character of the cyclic group.  The klt condition
makes its squared volume locally integrable, including after resolution,
so it extends holomorphically: a meromorphic top form with a divisorial
pole cannot be locally square integrable.
To see the orders directly, consider a boundary prime on the original
normal fiber with coefficient $u=p/e\in(0,1)$ in lowest terms.  At a
general point the cover has local coordinate $x=y^e$, and the pulled-back
root form has order
\begin{equation}
\label{adj:root-order}
                         e-1-eu=e-1-p\in\{0,\ldots,e-2\}.
\end{equation}
Away from the boundary its order at a nonexceptional prime is zero.
The corresponding calculation on a crepant resolution uses coefficients
less than one and proves integrability at exceptional divisors as well.

Suppose that $\tau'$ is another holomorphic top form of the same
character on a fiber.  The ratio $\tau'/\tau$ is an invariant meromorphic
function and descends to the original normal connected fiber.
A pole along a boundary prime downstairs would pull back with order at
least $e$, whereas \eqref{adj:root-order} permits only $e-2$ zeros of
$\tau$.  No such pole is possible.  Away from the boundary, $\tau$ has
no divisorial zeros downstairs.  Normality extends the ratio across
codimension two, and compactness makes it constant.  This also handles
a disconnected cover, since the full root group acts transitively on
its components over a connected fiber.  The highest Hodge piece of the
chosen character is therefore one-dimensional.

Finally, changing a local frame of $m b^*H$ changes the root by an
$m$th root of the corresponding base factor; changing the base volume
form contributes the inverse canonical transition.  Taking the $m$th tensor power removes
the finite root ambiguities.  These are precisely the transitions of
$m(b^*H-K_S)$, proving the asserted identity on $S^\circ$.
\end{proof}

\begin{lemma}[Extension across the discriminant]
\label{adj:extension}
The parabolic extension of the eigenline in Lemma~\ref{adj:root} is $L_S$
of \eqref{adj:discriminant}.  This identification persists on higher
smooth base models with simple normal crossing complement.
\end{lemma}

\begin{proof}
Fix a general point of a prime $Q_0=(t=0)$ on $S$.  We normalize the
root form by a nonvanishing local frame of $b^*H-K_S$.  More explicitly,
if the meromorphic section of $m b^*H$ chosen in Lemma~\ref{adj:root}
is $a$ times a nonvanishing local frame, divide its root by $a^{1/m}$
and by the local base volume form.  This is an ordinary frame after a
finite power substitution; the finite ambiguity is precisely part of
the parabolic convention.  Thus zeros or poles of the chosen meromorphic
base section do not enter the following exponent.
Tangential base coordinates may be treated as parameters.  On a log resolution,
\[
                         t=\prod_{i=1}^k x_i^{a_i}
\]
up to a nowhere-zero factor, and the squared absolute root form has
vertical density factors $|x_i|^{-2u_i}$.  Horizontal boundary factors
are integrable because their coefficients are less than one.
The squared Hodge norm of the relative form is its fiber integral.
Writing $s_i=-\log|x_i|$, the integral is computed on slices
\[
                 \sum_i a_i s_i=-\log|t|
\]
with exponential weight $\exp(-2\sum_i(1-u_i)s_i)$.
Conversion to the ordinary base area element contributes $|t|^{-2}$.
The minimum in \eqref{adj:threshold} thus gives constants $C,N>0$ such
that, on a sufficiently small punctured disk,
\begin{equation}
\label{adj:norm-power}
 C^{-1}|t|^{2(t_{Q_0}-1)}(-\log|t|)^{-N}
 \leq \|\tau\|^2
 \leq C|t|^{2(t_{Q_0}-1)}(-\log|t|)^N.
\end{equation}
Indeed, the slice has at most polynomial volume in $-\log|t|$, which
proves the upper estimate after factoring out the least exponential
decay.  For the lower estimate take a chart along a component attaining
the minimum, away from all other vertical components, and integrate
over a fixed compact set in the remaining directions.  Compactness
allows finitely many charts; the integrable horizontal factors do not
alter the power of $|t|$.

After a local power substitution making monodromy unipotent, extending
Hodge frames and their duals have norms bounded by powers of
$-\log|t|$.  These are the nilpotent-orbit estimates of
Schmid and Cattani--Kaplan--Schmid \cite{Schmid,CKS}, in the form used in
\cite[Section~4.1]{LI}.  Consequently the comparison between a
frame of $b^*H-K_S$ and an extending Hodge frame has rational order
$t_{Q_0}-1$.  Multiplication by $t^{1-t_{Q_0}}$, on a cover where this
power is integral, removes exactly that order.  Both the corrected
comparison and its inverse have at most logarithmic growth.  A
holomorphic function on a punctured disk with such growth has no pole;
applying this to the inverse excludes a zero.  The same argument with
tangential parameters extends the identification across each divisor,
and normal extension treats codimension two.  Thus the extending line is
\[
              b^*H-K_S+\sum_{Q_0}(t_{Q_0}-1)Q_0=L_S.
\]
For the last assertion, first make the boundary monodromies
unipotent.  In canonical extending frames a local monomial pullback
$t_i=\text{unit}\cdot\prod_j s_j^{a_{ij}}$ replaces the commuting
nilpotent residues $N_i$ by $\sum_i a_{ij}N_i$, which are again
nilpotent.  Both the canonical flat extension and its extending Hodge
filtration therefore pull back without an additional divisor.
Descending with rational weights incorporates the finite-monodromy
parts and gives the same conclusion for the parabolic line, including
exceptional base divisors; see \cite[Section~4]{LI}.
\end{proof}

We have now identified the moduli line with an extreme Hodge line, so
Lemma~\ref{adj:period-input} makes it nef and gives an actual line
descent.  To obtain a big line on the quotient, we must compare the
variation seen by this one line with that seen by the whole family.
Effectivity of the original boundary is decisive in this comparison.

\subsection{The line detects the full period rank}

\begin{lemma}[Equality of period ranks]
\label{adj:rank}
For the family of Lemma~\ref{adj:root}, the line period map of its root
eigenline and the full period map of $R^lh_*\R$ have the same generic
differential rank.
\end{lemma}

\begin{proof}
Work in an open patch on which the ranks are constant, and let $\xi$
be a holomorphic base vector field in the kernel of the line period
map.  Choose a local frame $\tau$ of the eigenline, viewed as a
relative top form on $V^\circ$.  The infinitesimal period formula says
that contraction with $\tau$ sends the Kodaira--Spencer class of $\xi$
to zero in the fiber cohomology $H^1(\Omega^{l-1})$.

Here is the corresponding lifting construction.  Pull back the tangent
sequence to the line generated by $\xi$ and push it out along
\[
 T_{V^\circ/S^\circ}\longrightarrow
 T_{V^\circ/S^\circ}\otimes K_{V^\circ/S^\circ}
       \simeq\Omega^{l-1}_{V^\circ/S^\circ},
 \qquad v\longmapsto\iota_v\tau.
\]
Write $U$ for a sufficiently small Stein base patch and put
$\mathcal F=\Omega^{l-1}_{V_U/U}$.  The pushed-out extension is
\[
             0\longrightarrow\mathcal F
               \longrightarrow\mathcal E
               \longrightarrow\cO_{V_U}\longrightarrow0,
\]
with class in $H^1(V_U,\mathcal F)$.  Shrink within the generic locus
so that $R^1h_*\mathcal F$ is locally free and coherent base change
holds.  The contraction formula makes the image of this class in
$H^0(U,R^1h_*\mathcal F)$ zero.  The Leray exact sequence and the
vanishing $H^1(U,h_*\mathcal F)=0$ show that
$H^1(V_U,\mathcal F)\to H^0(U,R^1h_*\mathcal F)$ is injective.
The extension class is therefore zero, and a holomorphic splitting
exists on the entire $V_U$.
Where $\tau$ is nonzero, divide the vertical part of the splitting by
$\tau$.  This gives a meromorphic vector field $\widetilde\xi$ on the
resolved cover projecting to $\xi$, with poles bounded by the zero
divisor of $\tau$.  Average it under the root group.  Its projection
remains $\xi$, and the averaged field is invariant.

We check this field at codimension-one points of the original normal
total space.  At a boundary prime use the ramification coordinate
$x=y^e$ from \eqref{adj:root-order}.  A tangential coefficient of an
invariant vector field has Laurent exponents congruent to zero modulo
$e$.  Its permitted pole order is at most $e-2$, so every exponent is
nonnegative.  The coefficient of $\partial/\partial y$ has exponents
congruent to one modulo $e$.  Its first possible negative exponent is
$1-e$, which is again excluded.  The normal coefficient is therefore
divisible by $y$.  Thus the descended vector field is holomorphic and
tangent to the boundary.  At a prime outside the boundary the root
form has no zero, and the conclusion is immediate.

The resolution is an isomorphism at the codimension-one points just
used.  Remaining exceptional divisors map into codimension at least
two in the original normal space.  Holomorphic derivations extend
across that subset, so the descended field is holomorphic everywhere
on the family over our base patch and still projects to $\xi$.
Its local flow preserves the boundary components, their coefficients,
and the singular locus.  By the functorial choice made above, the flow
first lifts to $X$ and preserves its crepant boundary $D_X$.

This flow then lifts holomorphically to the root cover.  Indeed, on the good
base patch the defining pluriform has divisor exactly the negative
multiple of the boundary.  Pullback by a pair-preserving flow gives a
pluriform with the same divisor.  Their ratio is a holomorphic unit on
the normal space; its $m$th root can be chosen starting at one in the
flow parameter.  The root normalized at flow time zero is unique as a germ.  Thus the
local cover lifts agree on overlaps of the unbranched open and glue
by normality.  Functorial
resolution then lifts the flow to $V^\circ$.
Consequently the Kodaira--Spencer class of the resolved family vanishes
on $\xi$, and so does its full period differential.
The reverse kernel inclusion holds because the line is part of the
full Hodge filtration.  The two kernels, and therefore the generic
ranks, are equal.
\end{proof}

The argument is the analytic counterpart of Ambro's infinitesimal
comparison \cite[Proposition~2.1]{Ambro}.  Notice exactly where it uses
$D\geq0$: the strict pole bound in \eqref{adj:root-order} is imposed at
primes of the original pair.  Negative crepant exceptional coefficients
are allowed because their images have codimension at least two.

\subsection{Replacing the moduli line by a klt boundary}

\begin{proof}[Proof of Proposition~\ref{adj:base}]
Apply Lemma~\ref{adj:period-input} to the variation of
Lemma~\ref{adj:root}.  Its monodromy is quasi-unipotent by the monodromy
theorem for this proper K\"ahler family \cite{Schmid}.  Lemma~\ref{adj:extension}
identifies its parabolic line with $L_S$.  Passing to the indicated
higher smooth projective model, and keeping the notation $S$, we have
\[
                         L_S\qsim p^*P
\]
for a nef rational line $P$ on a smooth projective quotient $Q$.
Let $r$ be the generic rank of the line period map.  Numerical dimension
of a nef line is unchanged by a surjective pullback, so
\[
 r=\nu(L_S)=\nu(P)\leq\dim Q
       \leq\operatorname{rank}(d\Phi)=r,
\]
where $\Phi$ is the full period map and the last equality is
Lemma~\ref{adj:rank}.  Thus $P$ is big if $\dim Q>0$.  If $Q$ is a point,
$L_S$ is rationally trivial.

Suppose first that $\dim Q>0$.  By Kodaira's lemma write
$P\qsim A+E$, with $A$ ample rational and $E\geq0$ rational.  For every
small positive rational $\varepsilon$,
\begin{equation}
\label{adj:small-fixed-part}
 P\qsim A_\varepsilon+\varepsilon E,
 \qquad A_\varepsilon=(1-\varepsilon)P+\varepsilon A
                 \quad\text{ample}.
\end{equation}
The pair $(S,\Delta_S)$ is sub-klt.  On a fixed log resolution of
$\Delta_S+p^*E$, choose $\varepsilon$ sufficiently small that
$(S,\Delta_S+\varepsilon p^*E)$ remains sub-klt.  Choose a sufficiently
large divisible integer $N$ and a general member
$G\in|N A_\varepsilon|$.  The system $|Np^*A_\varepsilon|$ is
basepoint-free; on that resolution its general member meets the fixed
boundary transversely.  Taking $N$ large makes its coefficient $1/N$
less than one.  Bertini therefore gives an effective rational divisor
\[
                    E_S=\varepsilon p^*E+\frac1N p^*G
                    \qsim L_S
\]
with $(S,\Delta_S+E_S)$ sub-klt.  In the point case take $E_S=0$.
In either case,
\begin{equation}
\label{adj:upstairs-presentation}
                 K_S+\Delta_S+E_S\qsim b^*H.
\end{equation}

Set $D_Z=b_*(\Delta_S+E_S)$.  This divisor is effective.  Indeed, for
each prime of the original normal base, some component of its pullback
on $Y$ dominates it.  That component has multiplicity $a\geq1$ and
boundary coefficient $u\geq0$, so its threshold is at most
$(1-u)/a\leq1$.  Thus the coefficient of the discriminant at every
nonexceptional base prime is nonnegative, and $E_S$ is effective.

Choose compatible canonical divisors and a rational Cartier divisor
representing $H$.  The line identity
\eqref{adj:upstairs-presentation} gives an integer $m>0$ and a rational
function $\varphi$ in the common function field of $S$ and $Z$ such that
\[
 K_S+\Delta_S+E_S-b^*H=\frac1m\operatorname{div}_S(\varphi).
\]
Pushing down gives
\[
 K_Z+D_Z-H=\frac1m\operatorname{div}_Z(\varphi).
\]
In particular, $K_Z+D_Z$ is rational Cartier and rationally linearly
equivalent to $H$.  Pulling the latter divisor equality back and
comparing with the former proves the crepant identity
\[
                 K_S+\Delta_S+E_S=b^*(K_Z+D_Z).
\]
Since the pair upstairs is sub-klt, the effective pair $(Z,D_Z)$ is klt.
This proves the proposition.
\end{proof}
